% Flavor catalog per-process template.
% process_id: EW001
% family: top_higgs_ew
% owner_agent_id: PKA-EW001
% writer_agent_id: null
% checker_agent_id: null
% opus_signoff_id: null
% Canonical metadata sidecar: EW001.yaml
% Status is controlled by the YAML sidecar status_history, not this comment.

\section{\texorpdfstring{\(S,T,U\) Oblique Electroweak Parameters}{S,T,U Oblique Electroweak Parameters}}

\subsection*{Process}
This entry covers the Peskin--Takeuchi oblique parameters \(S\), \(T\), and
\(U\), treated as universal new-physics corrections to electroweak gauge-boson
self energies.  The process is not flavor changing, but it is part of the
\texttt{top\_higgs\_ew} family because it is a leading precision bound on
warped, vector-like, composite, and custodial electroweak sectors.

\subsection*{PDG / equivalent value(s)}
The canonical values are the PDG 2025 electroweak-review fit in
\texttt{flavor\_catalog/references/EW001/pdg\_2025\_electroweak\_stu.txt},
accessed on 2026-05-16.  With \(U=0\), PDG quotes
\(S=0.026\pm0.075\), \(T=0.047\pm0.066\), and
\(\rho(S,T)=0.90\).  Letting all three float gives
\(S=0.021\pm0.096\), \(T=0.04\pm0.12\), and
\(U=0.008\pm0.092\), with correlations recorded in the sidecar.
For comparison, the public Gfitter oblique page gives
\(S=0.05\pm0.11\), \(T=0.09\pm0.13\), and \(U=0.01\pm0.11\).
The CDF-aware HEPfit/de Blas 2022 study quotes
\(M_W=80.4335\pm0.0094\,\mathrm{GeV}\); in its standard-average
\(U=0\) oblique fit, \(S=0.100\pm0.073\) and \(T=0.202\pm0.056\),
while its floating-\(U\) fit has \(U=0.134\pm0.087\).

\subsection*{Relevance to the RS / anarchic-flavor pipeline}
Bulk RS models generically shift \(S\) through mixing of the SM electroweak
bosons with KK gauge modes, and shift \(T\) when custodial symmetry is broken
by gauge, Higgs, or fermion embeddings.  PDG summarizes the warped contribution
as \(S\simeq 30 v^2/M_{\rm KK}^2\) and quotes \(S\leq0.18\) at 95\%
probability as implying \(M_{\rm KK}\gtrsim3.2\,\mathrm{TeV}\).  These bounds
are orthogonal to the catalog's \(\Delta F=2\) constraints, but they test the
same KK spectrum and localization choices before flavor observables are even
evaluated.

\subsection*{Post-2008 developments}
Relative to the CFW baseline, arXiv:0804.1954, the important change is not a
new LEP run but a much more mature global-fit environment: Higgs and top inputs
are measured directly, full two-loop and selected higher-order SM calculations
are included, and PDG now presents both \(U=0\) and floating-\(U\) oblique fits.
The 2022 CDF \(W\)-mass result created a distinct scenario in which \(T\) or
\(U\) can be pulled positive; the 2025 PDG review treats that impact as a
global-fit issue rather than replacing the main STU table with the CDF-only
preference.

\subsection*{Constraint validity and model dependence}
The STU language is cleanest for heavy, approximately universal new physics
whose leading effects are vacuum-polarization corrections.  It is less complete
when vertex corrections, non-universal fermion couplings, light new states, or
SMEFT flat directions are numerically important.  In RS applications, custodial
protection, brane kinetic terms, Higgs localization, and vector-like fermion
embeddings can change the translation from a fitted \(S,T,U\) point to a KK
mass bound.

\subsection*{Code coverage in this repo}
\textbf{NO}.  Greps over \texttt{quarkConstraints/}, \texttt{qcd/},
\texttt{flavorConstraints/}, \texttt{neutrinos/}, \texttt{yukawa/},
\texttt{warpConfig/}, \texttt{solvers/}, \texttt{scanParams/}, and
\texttt{tests/} found no S/T/U, oblique, EWPO, or W-mass-fit implementation.
Only generic electroweak constants and VEV comments surfaced, for example
\texttt{qcd/constants.py:11}, \texttt{qcd/running.py:3},
\texttt{quarkConstraints/qcd\_running.py:99}, \texttt{yukawa/neutrino.py:14},
and \texttt{yukawa/charged\_lepton.py:11}.

\subsection*{Implementation difficulty}
\textbf{HIGH}.  A production constraint would need a new electroweak precision
likelihood, correlation handling for the PDG/GFitter/HEPfit fit conventions,
and RS matching from the gauge and fermion KK spectra to vacuum-polarization
self energies or equivalent SMEFT operators.  This is a new electroweak
mode/matching calculation, not a reuse of the existing \(\Delta F=2\) basis.

\subsection*{Key references}
\texttt{pdg\_2025\_electroweak\_stu};
\texttt{gfitter\_oblique\_parameters};
\texttt{hepfit\_deblas\_2022\_cdf\_stu};
\texttt{hepfit\_precision\_ew\_page};
\texttt{peskin\_takeuchi\_1992\_oblique};
\texttt{cfw\_2008\_rs\_flavor}.
